\newcommand{\projectFullTitle}{\projectTitle : \projectSubtitle}
\newcommand{\projectTitle}{Battle of Languages}
\newcommand{\projectSubtitle}{A Usability Study of Functional Programming in Gameplay Programming}
\newcommand{\group}{PT103F19}
%%%%%%%%%%%%%%%%%%%%%%%%%%%%%%%%%%%%%%%%%%%%%%%%
% Language, Encoding and Fonts
% http://en.wikibooks.org/wiki/LaTeX/Internationalization
%%%%%%%%%%%%%%%%%%%%%%%%%%%%%%%%%%%%%%%%%%%%%%%%
% Select encoding of your inputs. Depends on
% your operating system and its default input
% encoding. Typically, you should use
%   Linux  : utf8 (most modern Linux distributions)
%            latin1
%   Windows: ansinew
%            latin1 (works in most cases)
%   Mac    : applemac
% Notice that you can manually change the input
% encoding of your files by selecting "save as"
% an select the desired input encoding.
\usepackage{scrextend}
\usepackage[utf8]{inputenc}
\usepackage[nointegrals]{wasysym}
%\usepackage{float}
\usepackage{xspace}
\usepackage{floatrow}
\usepackage{booktabs}
\usepackage{multicol}

\usepackage{siunitx}
\usepackage{csvsimple}
\usepackage[table,xcdraw]{xcolor}
% Make latex understand and use the typographic
% rules of the language used in the document.
\usepackage{array}
\usepackage[english]{babel}
% Use the vector font Latin Modern which is going
% to be the default font in latex in the future.
\usepackage{lmodern}
% Choose the font encoding
\usepackage[T1]{fontenc}
\usepackage{lastpage} % Tilføjet af Jesper, får sidetal til at virke.
%%%%%%%%%%%%%%%%%%%%%%%%%%%%%%%%%%%%%%%%%%%%%%%%
% Graphics and Tables
% http://en.wikibooks.org/wiki/LaTeX/Importing_Graphics
% http://en.wikibooks.org/wiki/LaTeX/Tables
% http://en.wikibooks.org/wiki/LaTeX/Colors
%%%%%%%%%%%%%%%%%%%%%%%%%%%%%%%%%%%%%%%%%%%%%%%%
% load a colour package
\usepackage[table]{xcolor}

\definecolor{aaublue}{RGB}{33,26,82}% dark blue
\definecolor{editgreen}{HTML}{008000}% dark green
% The standard graphics inclusion package
\usepackage{graphicx}
% Set up how figure and table captions are displayed
\usepackage[within=none]{caption}
\captionsetup{%
  font=footnotesize,% set font size to footnotesize
  labelfont=bf % bold label (e.g., Figure 3.2) font
}
% Make the standard latex tables look so much better
\usepackage{array,booktabs}
% Enable the use of frames around, e.g., theorems
% The framed package is used in the example environment
\usepackage{framed}

\usepackage{pgfplots}
\usepackage{pgfplotstable}
\pgfplotsset{compat=1.15}
\usepackage{tikz}
\usetikzlibrary{arrows,arrows.meta, fit, backgrounds, positioning, automata, shapes}
\usepackage{enumitem}
%%%%%%%%%%%%%%%%%%%%%%%%%%%%%%%%%%%%%%%%%%%%%%%%
% Mathematics
% http://en.wikibooks.org/wiki/LaTeX/Mathematics
%%%%%%%%%%%%%%%%%%%%%%%%%%%%%%%%%%%%%%%%%%%%%%%%
% Defines new environments such as equation,
% align and split
\usepackage{amsmath}
% Adds new math symbols
\usepackage{amssymb}
% Use theorems in your document
% The ntheorem package is also used for the example environment
% When using thmmarks, amsmath must be an option as well. Otherwise \eqref doesn't work anymore.
\usepackage[framed,amsmath,thmmarks]{ntheorem}

%%%%%%%%%%%%%%%%%%%%%%%%%%%%%%%%%%%%%%%%%%%%%%%%
% Page Layout
% http://en.wikibooks.org/wiki/LaTeX/Page_Layout
%%%%%%%%%%%%%%%%%%%%%%%%%%%%%%%%%%%%%%%%%%%%%%%%
% Change margins, papersize, etc of the document
\usepackage[top=1in, bottom=1.5in]{geometry}%, left=1in, right=1in]{geometry}

\setlength\parindent{0pt}   % Remove paragraph indentation
\setlength\parskip{1em}     % Sets paragraph spacing
%\usepackage[
%  inner=28mm,% left margin on an odd page
%  outer=41mm,% right margin on an odd page
%  ]{geometry}
% Modify how \chapter, \section, etc. look
% The titlesec package is very configureable

%%%%% Titles and sections
\usepackage{titlesec}
\newcommand{\hsp}{\hspace{20pt}}
\titleformat{\chapter}[hang]{\Huge\bfseries}{\thechapter\hsp|\hsp}{0pt}{}[]
\titleformat*{\section}{\normalfont\Large\bfseries}
\titleformat*{\subsection}{\normalfont\large\bfseries}
\titleformat*{\subsubsection}{\normalfont\normalsize\bfseries}

\titlespacing*{\chapter}{0pt}{0.5cm}{1cm}
\titlespacing*{name=\chapter,numberless}{0pt}{0.5cm}{1cm}

% Clear empty pages between chapters
\let\origdoublepage\cleardoublepage
\newcommand{\clearemptydoublepage}{%
  \clearpage
  {\pagestyle{empty}\origdoublepage}%
}
\let\cleardoublepage\clearemptydoublepage

% Change the headers and footers
\usepackage{fancyhdr}
\pagestyle{fancy}
\fancyhf{} %delete everything
\renewcommand{\headrulewidth}{0pt} %remove the horizontal line in the header
\fancyfoot[C]{\thepage} %page number on all pages
\setlength{\headheight}{14pt}
\setlength{\marginparwidth}{2cm}
% Do not stretch the content of a page. Instead,
% insert white space at the bottom of the page
\raggedbottom
% Enable arithmetics with length. Useful when
% typesetting the layout.
\usepackage{calc}

\usepackage[
%  disable, %turn off notes
  colorinlistoftodos, %enable a coloured square in the list of todos
  textwidth=\marginparwidth, %set the width of the todonotes
  textsize=scriptsize, %size of the text in the todonotes
  ]{todonotes}

\usepackage[T1]{fontenc}
\usepackage{inconsolata}
\usepackage{color}
\usepackage{svg}

\usepackage{listings}
% see, e.g., http://en.wikibooks.org/wiki/LaTeX/Customizing_LaTeX#New_commands
% for more information on how to create macros

%%%%%%%%%%%%%%%%%%%%%%%%%%%%%%%%%%%%%%%%%%%%%%%%
% Macros for the titlepage
%%%%%%%%%%%%%%%%%%%%%%%%%%%%%%%%%%%%%%%%%%%%%%%%
%Creates the aau titlepage
\newcommand{\aautitlepage}[3]{%
  {
    %set up various length
    \ifx\titlepageleftcolumnwidth\undefined
      \newlength{\titlepageleftcolumnwidth}
      \newlength{\titlepagerightcolumnwidth}
    \fi
    \setlength{\titlepageleftcolumnwidth}{0.5\textwidth-\tabcolsep}
    \setlength{\titlepagerightcolumnwidth}{\textwidth-2\tabcolsep-\titlepageleftcolumnwidth}
    %create title page
    \thispagestyle{empty}
    \noindent%
    \begin{tabular}{@{}ll@{}}
      \parbox{\titlepageleftcolumnwidth}{
        \iflanguage{danish}{%
          \includegraphics[width=\titlepageleftcolumnwidth]{00-images/aau-logo-da}
        }{%
          \includegraphics[width=\titlepageleftcolumnwidth]{00-images/aau-logo-en}
        }
      } &
      \parbox{\titlepagerightcolumnwidth}{\raggedleft\sf\small
        #2
      }\bigskip\\
       #1 &
      \parbox[t]{\titlepagerightcolumnwidth}{%
      \textbf{Abstract:}\bigskip\par
        \fbox{\parbox{\titlepagerightcolumnwidth-2\fboxsep-2\fboxrule}{%
          #3
        }}
      }\\
    \end{tabular}
    \vfill
    \iflanguage{danish}{%
      \noindent{\footnotesize\emph{Rapportens indhold er frit tilgængeligt, men offentliggørelse (med kildeangivelse) må kun ske efter aftale med forfatterne.}}
    }{%
      \noindent{\footnotesize\emph{The content of this report is freely available, but publication (with reference) may only be pursued due to agreement with the author.}}
    }
    \clearpage
  }
}

%Create english project info
\newcommand{\englishprojectinfo}[8]{%
  \parbox[t]{\titlepageleftcolumnwidth}{
    \textbf{Title:}\\ #1\bigskip\par
    \textbf{Theme:}\\ #2\bigskip\par
    \textbf{Project Period:}\\ #3\bigskip\par
    \textbf{Project Group:}\\ #4\bigskip\par
    \textbf{Participant(s):}\\ #5\bigskip\par
    \textbf{Supervisor(s):}\\ #6\bigskip\par
    \textbf{Copies:} #7\bigskip\par
    \textbf{Page Numbers:} \pageref{LastPage}\bigskip\par
    \textbf{Date of Completion:}\\ #8
  }
}

%Create danish project info
\newcommand{\danishprojectinfo}[8]{%
  \parbox[t]{\titlepageleftcolumnwidth}{
    \textbf{Titel:}\\ #1\bigskip\par
    \textbf{Tema:}\\ #2\bigskip\par
    \textbf{Projektperiode:}\\ #3\bigskip\par
    \textbf{Projektgruppe:}\\ #4\bigskip\par
    \textbf{Deltagere:}\\ #5\bigskip\par
    \textbf{Vejleder:}\\ #6\bigskip\par
    \textbf{Oplagstal:} #7\bigskip\par
    \textbf{Sidetal:} \pageref{LastPage}\bigskip\par
    \textbf{Afleveringsdato:}\\ #8
  }
}

\usepackage{courier}
\definecolor{gray}{rgb}{0.4,0.4,0.4}
\definecolor{darkblue}{rgb}{0.0,0.0,0.6}
\definecolor{cyan}{rgb}{0.0,0.6,0.6}


%\setmonofont{Consolas} %to be used with XeLaTeX or LuaLaTeX
\definecolor{bluekeywords}{rgb}{0,0,1}
\definecolor{greencomments}{rgb}{0,0.5,0}
\definecolor{redstrings}{rgb}{0.64,0.08,0.08}
\definecolor{xmlcomments}{rgb}{0.5,0.5,0.5}
\definecolor{types}{rgb}{0.17,0.57,0.68}

% PDF package
\usepackage[final]{pdfpages}
\usepackage{glossaries}
\usepackage{longtable}
\usepackage{tasks}
\usepackage{subcaption}
\usepackage{setspace}
\addtokomafont{labelinglabel}{\sffamily\bfseries}
\usepackage{dirtytalk}
\usepackage{etoolbox}
\usepackage{silence}

\usepackage[backend=biber, style=ieee, urldate=long]{biblatex} % biblatex - bibliography tool %
\bibliography{00-preliminary/bib}
\renewcommand*{\lstlistlistingname}{List of Listings}

\setlength{\headheight}{15pt}

\usepackage{rotating}
\usepackage{tablefootnote}

% Style to select only points from #1 to #2 (inclusive)
\pgfplotsset{select coords between index/.style 2 args={
    x filter/.code={
        \ifnum\coordindex<#1\def\pgfmathresult{}\fi
        \ifnum\coordindex>#2\def\pgfmathresult{}\fi
    }
}}

\usepackage[outputdir=build]{minted}
%%%%%%%%%% C# %%%%%%%%%%
\setminted[csharp]{
frame=single,%lines
framesep=2mm,
fontsize=\footnotesize,
linenos,
breaklines
}
\usemintedstyle[csharp]{vs}

\setminted[fsharp]{
frame=single,%lines
framesep=2mm,
fontsize=\footnotesize,
linenos,
breaklines
}
\usemintedstyle[fsharp]{vs}

% Quotations package
\usepackage[autostyle]{csquotes}

%%%%%%%%%%%%%%%%%%% Packages %%%%%%%%%%%%%%%%%%%
\usepackage{xifthen}
\usepackage{xparse}


%%%%%%%%%%%%%%%%%%% PGF Styles%%%%%%%%%%%%%%%%%%%
\pgfplotsset{
  abs log x ticks/.style={
      xticklabel={
        \pgfkeys{/pgf/fpu=true}
        \pgfmathparse{exp(\tick)}%
        \pgfmathprintnumber[fixed relative, precision=3]{\pgfmathresult}
        \pgfkeys{/pgf/fpu=false}
      }
  },
  abs log y ticks/.style={
      yticklabel={
        \pgfkeys{/pgf/fpu=true}
        \pgfmathparse{exp(\tick)}%
        \pgfmathprintnumber[fixed relative, precision=3]{\pgfmathresult}
        \pgfkeys{/pgf/fpu=false}
      }
  }
}

%%%%%%%%%%%%%%%%%%% Commands %%%%%%%%%%%%%%%%%%%
\newcommand{\ifequals}[3]{\ifthenelse{\equal{#1}{#2}}{#3}{}}
\newcommand{\case}[2]{#1 #2} % Dummy, so \renewcommand has something to overwrite...
\newenvironment{switch}[1]{\renewcommand{\case}{\ifequals{#1}}}{}

\newcommand{\symbolicsStrat}{%
symbolic x coords={,
    Problem,
    10,
    100,
    1000,
    10000,
    100000,
}
}

\newcommand{\symbolic}[1]{symbolic x coords={#1}}
\NewDocumentCommand\lnm{O{y}O{2}}{%
\begin{switch}{#1}%
    \case{x}{xmode=log, log basis x={#2}}%
    \case{y}{ymode=log, log basis y={#2}}%
\end{switch}%
}
\newcommand{\plotData}[2]{
    \addplot table [y={#1}] {#2};
    \addlegendentry{#1};
}

\newcommand{\plotUnmarkedData}[2]{
  \addplot+[mark=none] table [y={#1}] {#2};
  \addlegendentry{#1};
}

\NewDocumentCommand\optPlotData{O{}mm}{
  \addplot+[#1] table [y={#2}] {#3};
  \addlegendentry{#2};
}

\newcommand{\plotDataWithError}[2]{
    \addplot+[error bars/.cd, y dir=both,y explicit] table [y={#1}, y error={#1 Error}] {#2};
    \addlegendentry{#1};
}

\newcommand{\plotDataWithErrorAndLegend}[3]{
    \addplot+[error bars/.cd, y dir=both,y explicit] table [y={#1}, y error={#1 Error}] {#2};
    \addlegendentry{#3};
}

\newcommand{\logmode}[1]{\IfBooleanTF{#1}{log}{normal}}
\newcommand{\logbase}[1]{\IfBooleanTF{#1}{2}{1}}

% O/o for optional, M/m for mandatory, s for star
% s #1 BooleanTrue if star, BooleanFalse otherwise
% O #2 PGFPlot options
% O #3 Y label
% O #4 X label
% m #5 Caption
% m #6 Label
% m #7 Plots
\NewDocumentCommand\lineChart{sO{}O{Logarithmic Run Time (ns)}O{Lower is better}mmm}{%

\begin{figure}[H]
\makebox[\textwidth][c]{
\tikzsetnextfilename{#6}
\begin{tikzpicture}
    \begin{axis}[
        ylabel={#3},
        xlabel={#4},
        width=1.2\textwidth,
        height=8cm,
        ymajorgrids,
        xtick=data,
        #2,
        ymode/.expand once=\logmode{#1},
        log basis y/.expand once=\logbase{#1},
        xticklabel style={rotate=45, anchor=east},
        legend columns = -1,
        legend style={
            draw=none,
            at={(0.5,1.05)},
            anchor=south,
            column sep=1ex,
            font=\tiny
        }
    ]
    #7
    \end{axis}
\end{tikzpicture}}
\caption{#5}
\label{fig:#6}
\end{figure}
}

% O/o for optional, M/m for mandatory, s for star
% s #1 BooleanTrue if star, BooleanFalse otherwise
% O #2 bar width
% O #3 PGFPlot options
% O #4 Y label
% O #5 X label
% m #6 Caption
% m #7 Label
% m #8 Plots
\NewDocumentCommand\barChart{sO{12}O{}O{Run Time (ns)}O{Lower is better}mmm}{%
    \begin{figure}[H]
    \makebox[\textwidth][c]{
    \tikzsetnextfilename{#7}
    \begin{tikzpicture}
        \begin{axis}[
            ylabel={#4},
            xlabel={#5},
            width=1.2\textwidth,
            height=8cm,
            ymajorgrids,
            xtick=data,
            ybar=0.4pt,
            bar width=#2pt,
            xticklabel style={rotate=45, anchor=east},
            #3,
            ymode/.expand once=\logmode{#1},
            log basis y/.expand once=\logbase{#1},
            legend columns = -1,
            legend style={
                draw=none,
                at={(0.5,1.05)},
                anchor=south,
                column sep=1ex,
                font=\tiny
            }
        ]
        #8
        \end{axis}
    \end{tikzpicture}}
    \caption{#6}
    \label{fig:#7}
\end{figure}
}


\usepackage{hyperref}
\hypersetup{%
	plainpages=false,%
	pdfauthor={\group\ - Tobias Morell, Thomas Gwynfryn McCollin},%
	pdftitle={\group\ - \projectTitle},%
	pdfsubject={Project Report},%
	bookmarksnumbered=true,%
	colorlinks,%
	citecolor=aaublue,%
	filecolor=aaublue,%
	linkcolor=aaublue,% you should probably change this to black before printing
	urlcolor=aaublue,%
	pdfstartview=FitH%
}

\pgfplotstableread[col sep = comma]{00-data/fsharp/accumulation.csv}\accumulationData
\pgfplotstableread[col sep = comma]{00-data/fsharp/sum.csv}\summationData
\pgfplotstableread[col sep = comma]{00-data/csharp/accumulation.csv}\accumulationDataCsharp
\pgfplotstableread[col sep = comma]{00-data/csharp/sum.csv}\summationDataCsharp

\pgfplotstableread[col sep = comma]{00-data/fsharp/matrix.csv}\matrixSumData
\pgfplotstableread[col sep = comma]{00-data/csharp/matrix.csv}\matrixSumDataCsharp

\pgfplotstableread[col sep = comma]{00-data/work-bias-wait-after-join.csv}\workWithDependencyData
\pgfplotstableread[col sep = comma]{00-data/work-bias-wait-before-join.csv}\workWithoutDependencyData
\pgfplotstableread[col sep = comma]{00-data/nary-work-bias-wait-after-join.csv}\workWithDependencyDataNary
\pgfplotstableread[col sep = comma]{00-data/nary-work-bias-wait-before-join.csv}\workWithoutDependencyDataNary

\pgfplotstableread[col sep = comma]{00-data/ai-benchmark-runtimes.csv}\aiBenchmarkRuntimesData
\pgfplotstableread[col sep = comma]{00-data/ai-benchmark.csv}\aiBenchmarkData
\pgfplotstableread[col sep = comma]{00-data/csharp-vs-frp.csv}\aiBenchmarkOverheadData
\pgfplotstableread[col sep = comma]{00-data/csharp-incremental-vs-normal.csv}\aiBenchmarkCsharpGCData
\pgfplotstableread[col sep = comma]{00-data/csharp-incremental-vs-normal-collapsed.csv}\aiCsharpGCDataCol
\pgfplotstableread[col sep = comma]{00-data/frp-incremental-vs-normal.csv}\aiBenchmarkFsharpGCData
\pgfplotstableread[col sep = comma]{00-data/frp-incremental-vs-normal-collapsed.csv}\aiFsharpGCDataCol
